\chapter{Estimation}

\section{Random Sample}

A sample $(Y_1, Y_2, \dots, Y_T)$ is called \textbf{i.i.d.} if
\[
f_T(y_1,\dots,y_T) = \prod_{t=1}^T f(y_t)
\]

\section{Estimators}

An estimator of $\theta \in \mathbb{R}^K$ is
\[
\hat{\theta} = \hat{\theta}(Y_1,\dots,Y_T)
\]

An estimate is
\[
\hat{\theta}(y_1,\dots,y_T)
\]

\section{Method of Moments}

Sample moment:
\[
\hat{\mu}_r = \frac{1}{T} \sum_{t=1}^T Y_t^r
\]

Estimator:
\[
\hat{\mu}_r = h_r(\hat \theta)
\]

\subsection{Example}

\[
\hat{\beta}_m = \frac{1}{T} \sum_{t=1}^T Y_t = \hat{\mu}_1
\]

\[
\begin{aligned}
\hat{\sigma}_m^2 
&= \frac{1}{T} \sum_{t=1}^T (Y_t - \hat{\beta}_m)^2 \\
&= \hat{\mu}_2 - \hat{\mu}_1^2
\end{aligned}
\]

\section{Maximum Likelihood}

Likelihood:
\[
\begin{aligned}
L(\theta) 
&= f_T(Y_1, \dots ,Y_t|\theta)\\
&= \prod_{t=1}^T f(Y_t|\theta) \qquad (i.i.d.)
\end{aligned}
\]

Log-likelihood:
\[
\ell(\theta) = \sum_{t=1}^T \ln f(Y_t|\theta)
\]

Estimator:
\[
\begin{aligned}
\hat{\theta}_{ML} 
&= \arg\max_{\theta \in \Omega} L(\theta)\\
&= \arg\max_{\theta \in \Omega} \ell(\theta)
\end{aligned}
\]

\subsection{Normal Example}

Let $(Y_1,\dots,Y_T)$ be an i.i.d. sample from
\[
Y_t \sim N(\beta,\sigma^2),
\qquad \theta=(\beta,\sigma^2),\ \sigma^2>0.
\]

The density is
\[
f(y_t|\beta,\sigma^2)
=
\frac{1}{\sqrt{2\pi\sigma^2}}
\exp\!\left(-\frac{(y_t-\beta)^2}{2\sigma^2}\right).
\]

Hence the likelihood is
\[
L(\beta,\sigma^2)
=
\prod_{t=1}^T f(Y_t|\beta,\sigma^2),
\]
and the log-likelihood is
\[
\ell(\beta,\sigma^2)
=
-\frac{T}{2}\ln(2\pi)
-\frac{T}{2}\ln(\sigma^2)
-\frac{1}{2\sigma^2}\sum_{t=1}^T (Y_t-\beta)^2.
\]

Maximizing $\ell$ yields the first-order conditions:
\[
\frac{\partial \ell}{\partial \beta}
=
\frac{\sum_{t=1}^T Y_t - T\beta}{\sigma^2}=0,
\]
\[
\frac{\partial \ell}{\partial \sigma^2}
=
-\frac{T}{2\sigma^2}
+
\frac{\sum_{t=1}^T (Y_t-\beta)^2}{2\sigma^4}=0.
\]

Solving gives the MLE:
\[
\hat{\beta}
=
\frac{1}{T}\sum_{t=1}^T Y_t,
\qquad
\hat{\sigma}^2
=
\frac{1}{T}\sum_{t=1}^T (Y_t-\hat{\beta})^2.
\]

\paragraph{Summary.}
The log-likelihood is concave in $(\beta,\sigma^2)$, so the first-order conditions characterize the global maximum. 
For the normal distribution, the MLE coincides with the method of moments:
the sample mean and (biased) sample variance.

\paragraph{Note:} 
The maximum likelihood method requires knowing the probability distribution of the Yt’s. (Other methods can be “less demanding”...)

\section{Least Squares}

Assume that the conditional mean of $Y_t$ is given by
\[
E(Y_t) = h_t(\theta),
\qquad t=1,\dots,T,
\qquad \theta \in \Omega \subset \mathbb{R}^K.
\]

Define the error term
\[
e_t = Y_t - h_t(\theta),
\]
so that $e_t$ is a random variable with mean zero:
\[
E(e_t) = E(Y_t) - h_t(\theta) = 0.
\]

The error sum of squares (ESS) is defined as
\[
S(\theta)
=
\sum_{t=1}^T (Y_t - h_t(\theta))^2.
\]

The least squares estimator is
\[
\hat{\theta}_{LS}
=
\arg\min_{\theta \in \Omega} S(\theta).
\]

\subsection{Example}

Let $Y_1,\dots,Y_T$ be a sample such that
\[
E(Y_t) = \beta,
\]
and hence $h_t(\beta)=\beta$.

Then the objective function becomes
\[
S(\beta)
=
\sum_{t=1}^T (Y_t - \beta)^2.
\]

Since $S(\beta)$ is convex in $\beta$, the first-order condition is sufficient for a minimum:
\[
\frac{\partial S}{\partial \beta}
=
\sum_{t=1}^T -2(Y_t - \beta)
=
0.
\]

Solving,
\[
\hat{\beta}
=
\frac{1}{T}\sum_{t=1}^T Y_t.
\]

\paragraph{Summary.}
The least squares estimator minimizes the sum of squared deviations from the conditional mean function.
In this simple case, it coincides with the sample mean.

\paragraph{Note.}
When $E(Y_t)=\beta$, the least squares, method of moments, and maximum likelihood estimators are identical:
\[
\hat{\beta}_{LS}
=
\hat{\beta}_{MM}
=
\hat{\beta}_{ML}.
\]

\section{Finite Sample Properties}

\subsection{Unbiasedness}

\[
E(\hat{\theta}) = \theta
\]

\subsection{Efficiency}

Smallest variance among unbiased estimators.

\subsection{Cram\'er--Rao Lower Bound}

Let $\hat{\theta}$ be an unbiased estimator of the parameter $\theta \in \mathbb{R}^K$.
Then, under regularity conditions,
\[
\operatorname{Var}(\hat{\theta}) \;\succeq\; I(\theta)^{-1},
\]
where
\[
I(\theta)
=
- E\!\left[\frac{\partial^2 \ell(\theta)}{\partial \theta \,\partial \theta'}\right]
\]
is the \emph{information matrix}, and $\ell(\theta)$ is the log-likelihood function.

\paragraph{Interpretation.}
The matrix $I(\theta)^{-1}$ is called the \emph{Cram\'er--Rao lower bound}.
It provides a lower bound on the variance--covariance matrix of any unbiased estimator of $\theta$, i.e.
\[
\operatorname{Var}(\hat{\theta}) - I(\theta)^{-1}
\]
is positive semidefinite.

\paragraph{Efficient estimator.}
An unbiased estimator $\hat{\theta}$ is said to be \emph{efficient} if
\[
\operatorname{Var}(\hat{\theta}) = I(\theta)^{-1}.
\]
In this case, no other unbiased estimator has a smaller variance.

\paragraph{Scalar case.}
If $\theta$ is scalar, then
\[
\operatorname{Var}(\hat{\theta}) \ge \frac{1}{I(\theta)}.
\]

\paragraph{Note.}
The information matrix depends on the log-likelihood function, and hence on the full probability distribution of the data.
Therefore, applying the Cram\'er--Rao bound requires knowledge of the distribution of $Y_t$.

\subsection{Best Linear Unbiased Estimator (BLUE)}

An estimator $\hat{\theta}$ is said to be \emph{best linear unbiased} (BLUE) if:

\begin{itemize}
\item It is \emph{linear} in the data:
\[
\hat{\theta} = \sum_{t=1}^T a_t Y_t
\]
for some constants $a_t$.

\item It is \emph{unbiased}:
\[
E(\hat{\theta}) = \theta.
\]

\item It has the \emph{smallest variance} among all linear unbiased estimators.
\end{itemize}

\paragraph{Note.}
The BLUE property does not require knowledge of the full probability distribution of $Y_t$.

\subsubsection{Example 1: Estimation of the Mean}

Let $Y_1,\dots,Y_T$ be a random sample with
\[
E(Y_t) = \beta,
\qquad
V(Y_t) = \sigma^2.
\]

Consider the estimator
\[
\hat{\beta}
=
\frac{1}{T}\sum_{t=1}^T Y_t.
\]

Then
\[
E(\hat{\beta})
=
\frac{1}{T}\sum_{t=1}^T E(Y_t)
=
\beta,
\]
so it is unbiased.

Its variance is
\[
V(\hat{\beta})
=
V\!\left(\frac{1}{T}\sum_{t=1}^T Y_t\right)
=
\frac{1}{T^2}\sum_{t=1}^T V(Y_t)
=
\frac{\sigma^2}{T}.
\]

Since $\hat{\beta}$ is linear and has the smallest variance among all linear unbiased estimators,
it is the \emph{BLUE} of $\beta$.


\paragraph{Connection with Cram\'er--Rao Bound.}


If, in addition, $Y_t \sim N(\beta,\sigma^2)$, then
\[
V(\hat{\beta}) = \frac{\sigma^2}{T} = I(\theta)^{-1}.
\]
Hence $\hat{\beta}$ is \emph{efficient}.

\paragraph{Conclusion.}
In this case,
\[
\hat{\beta}_{LS}
=
\hat{\beta}_{MM}
=
\hat{\beta}_{ML}
=
\frac{1}{T}\sum_{t=1}^T Y_t,
\]
so all methods coincide.
\subsection{Variance Estimator}

Let $Y_1,\dots,Y_T$ be a random sample with
\[
E(Y_t)=\beta,
\qquad
V(Y_t)=\sigma^2.
\]

Consider the estimator
\[
\hat{\sigma}^2
=
\frac{1}{T}\sum_{t=1}^T (Y_t - \hat \beta)^2,
\qquad
\hat \beta = \frac{1}{T}\sum_{t=1}^T Y_t.
\]

Using the decomposition
\[
Y_t - \hat \beta
=
(Y_t - \beta) - (\hat \beta - \beta),
\]
we obtain
\[
E(\hat{\sigma}^2)
=
E\!\left[
\frac{1}{T}\sum_{t=1}^T (Y_t-\beta)^2
-
(\hat \beta-\beta)^2
\right].
\]

Hence
\[
E(\hat{\sigma}^2)
=
\sigma^2 - V(\hat \beta)
=
\sigma^2 - \frac{\sigma^2}{T}
=
\frac{T-1}{T}\sigma^2
<
\sigma^2,
\]
so $\hat{\sigma}^2$ is a \emph{biased} estimator.

\paragraph{Unbiased estimator.}
Define
\[
\hat{\sigma}^2_u
=
\frac{1}{T-1}\sum_{t=1}^T (Y_t - \hat \beta)^2.
\]

Since
\[
\hat{\sigma}^2_u
=
\hat{\sigma}^2 \cdot \frac{T}{T-1},
\]
it follows that
\[
E(\hat{\sigma}^2_u)
=
\sigma^2,
\]
so $\hat{\sigma}^2_u$ is unbiased.

\paragraph{Summary.}
\[
\text{MLE / MOM: } \frac{1}{T}\sum (Y_t-\hat \beta)^2 \quad (\text{biased})
\]
\[
\text{Unbiased: } \frac{1}{T-1}\sum (Y_t-\hat \beta)^2
\]


\section{Asymptotic Properties}

\subsection{Consistency}

An estimator $\hat{\theta}$ of $\theta$ is said to be \emph{consistent} if, for any $\varepsilon > 0$,
\[
\lim_{T \to \infty} P\big(|\hat{\theta} - \theta| < \varepsilon \big) = 1.
\]

Equivalently, $\hat{\theta}$ is consistent if it converges in probability to $\theta$:
\[
\text{plim } \hat{\theta} = \theta.
\]

\paragraph{Sufficient conditions.}
A sufficient condition for consistency is:
\[
\lim_{T \to \infty} E(\hat{\theta}) = \theta
\quad (\text{asymptotic unbiasedness}),
\]
\[
\lim_{T \to \infty} \operatorname{Var}(\hat{\theta}) = 0.
\]

\paragraph{Example.}
Let $Y_1,\dots,Y_T$ be a random sample with
\[
E(Y_t) = \beta,
\qquad
V(Y_t) = \sigma^2.
\]

Consider the estimator
\[
\hat{\beta}
=
\frac{1}{T}\sum_{t=1}^T Y_t.
\]

Then
\[
E(\hat{\beta}) = \beta,
\qquad
V(\hat{\beta}) = \frac{\sigma^2}{T} \to 0.
\]

Hence,
\[
\text{plim } \hat{\beta} = \beta,
\]
so $\hat{\beta}$ is a consistent estimator of $\beta$.
\subsection{Central Limit Theorem}

Let $Y_1,\dots,Y_T$ be an i.i.d. sample with
\[
E(Y_t)=\beta,
\qquad
V(Y_t)=\sigma^2.
\]

For the sample mean
\[
\bar{Y} = \frac{1}{T}\sum_{t=1}^T Y_t,
\]
we have
\[
\sqrt{T}(\bar{Y} - \beta)
\xrightarrow{d}
N(0,\sigma^2).
\]

\paragraph{Implication.}
Equivalently, for large $T$,
\[
\bar{Y} \approx N\!\left(\beta, \frac{\sigma^2}{T}\right),
\]
regardless of the underlying distribution of $Y_t$.

\subsection{Asymptotic Normality}

More generally, an estimator $\hat{\theta}$ is said to be asymptotically normal if
\[
\sqrt{T}(\hat{\theta} - \theta)
\xrightarrow{d}
N(0,\Omega),
\]
for some variance matrix $\Omega$.

Equivalently,
\[
\hat{\theta}
\approx
N\!\left(\theta, \frac{\Omega}{T}\right)
\quad \text{for large } T.
\]

\subsection{Asymptotic Efficiency}

An estimator $\hat{\theta}$ is asymptotically efficient if it is consistent and
\[
\sqrt{T}(\hat{\theta} - \theta)
\xrightarrow{d}
N\!\left(0,\; \lim_{T\to\infty}\left[\frac{1}{T}I(\theta)\right]^{-1}\right),
\]
where
\[
I(\theta) = -E\!\left[\frac{\partial^2 \ell(\theta)}{\partial \theta \,\partial \theta'}\right]
\]
is the information matrix.

\paragraph{MLE property.}
Under regularity conditions, the maximum likelihood estimator $\hat{\theta}_{ML}$ is:
\begin{itemize}
\item consistent,
\item asymptotically normal,
\item asymptotically unbiased,
\item asymptotically efficient.
\end{itemize}

\subsection{Example (Normal Case)}

Let $Y_t \sim N(\beta,\sigma^2)$. Then the MLE
\[
\hat{\theta}_{ML} = (\hat{\beta}, \hat{\sigma}^2)
\]
satisfies
\[
\sqrt{T}(\hat{\theta}_{ML} - \theta)
\xrightarrow{d}
N\!\left(0,
\begin{pmatrix}
\sigma^2 & 0 \\
0 & 2\sigma^4
\end{pmatrix}
\right).
\]

Hence, asymptotically,
\[
V(\hat{\beta}) \approx \frac{\sigma^2}{T},
\qquad
V(\hat{\sigma}^2) \approx \frac{2\sigma^4}{T}.
\]